% GENERATED FILE. Edit canonical lessons, then run npm run generate:books.
% canonical-source-hash: fnv1a64:7011225eb75aa0d2
% canonical-lessons: PA-C79-sapp, PA-C79-vadda, PA-C79-chhota, PA-C79-navam, PA-C79-purana

\chapter{Snake, Big, Small, New, Old}
\label{ch:pa-snake-big-small-new-old}
\hlchaptermodality{79}

\begin{chapteropening}
\textbf{By the end of this chapter:} \emph{I can say a snake, big, small, new and old.}

Complete the last lesson of chapter 79: I can say a snake, big, small, new and old.
\end{chapteropening}

\section[\texorpdfstring{\pa{ਸੱਪ} (sapp)}{ਸੱਪ (sapp)}]{\pa{ਸੱਪ} (sapp) — a snake}

\label{lesson:PA-C79-sapp}



\begin{quote}
Before the new one: say the Punjabi for an elephant, then the Punjabi for a monkey.
\end{quote}

\subsection*{You'll want to know: \pa{ਸੱਪ}}
\textbf{\pa{ਸੱਪ}} — \emph{sapp} — \textquotedblleft{}a snake\textquotedblright{}.

\begin{cousinweb}
\textbf{\pa{ਸੱਪ}} (\emph{sapp}) is a tadbhava, a word inherited from Sanskrit and reshaped by sound change. It comes from \emph{sarpa}, \textquotedblleft{}snake\textquotedblright{} (CDIAL 13271).
\end{cousinweb}

\begin{grammarlens}[title={What you've built}]
One more word for this chapter.
\end{grammarlens}

\subsection*{Your turn}
\begin{itemize}
\raggedright
  \item \emph{Say it:} \emph{sapp}
  \item \emph{Say it:} \emph{sapp}, once more
  \item \emph{Say it:} say \emph{sapp}
  \item \emph{Recall:} say the Punjabi for an elephant, then the Punjabi for a monkey, then say \emph{sapp} again
\end{itemize}

\begin{tcolorbox}[breakable,colback=teal!4,colframe=teal!35!black,title={Before you move on}]
What does \pa{ਸੱਪ} mean? (\textquotedblleft{}A snake\textquotedblright{}.) Say it once more.
\end{tcolorbox}
\section[\texorpdfstring{\pa{ਵੱਡਾ} (vaḍḍā)}{ਵੱਡਾ (vaḍḍā)}]{\pa{ਵੱਡਾ} (vaḍḍā) — big}

\label{lesson:PA-C79-vadda}



\begin{quote}
Before the new one: say the Punjabi for a monkey, then the Punjabi for a snake.
\end{quote}

\subsection*{You'll want to know: \pa{ਵੱਡਾ}}
\textbf{\pa{ਵੱਡਾ}} — \emph{vaḍḍā} — \textquotedblleft{}big\textquotedblright{}.

\begin{cousinweb}
\textbf{\pa{ਵੱਡਾ}} (\emph{vaḍḍā}) is a tadbhava, a word inherited from Sanskrit and reshaped by sound change. It comes from \emph{vaḍra}, \textquotedblleft{}big\textquotedblright{} (CDIAL 11225). CDIAL notes that the Sanskrit word is known only from the old word-lists, and doubts a link with \emph{vṛddha}, \textquotedblleft{}grown\textquotedblright{}.
\end{cousinweb}

\begin{grammarlens}[title={What you've built}]
One more word for this chapter.
\end{grammarlens}

\subsection*{Your turn}
\begin{itemize}
\raggedright
  \item \emph{Say it:} \emph{vaḍḍā}
  \item \emph{Say it:} \emph{vaḍḍā}, once more
  \item \emph{Say it:} say \emph{vaḍḍā}
  \item \emph{Recall:} say the Punjabi for a monkey, then the Punjabi for a snake, then say \emph{vaḍḍā} again
\end{itemize}

\begin{tcolorbox}[breakable,colback=teal!4,colframe=teal!35!black,title={Before you move on}]
What does \pa{ਵੱਡਾ} mean? (\textquotedblleft{}Big\textquotedblright{}.) Say it once more.
\end{tcolorbox}
\section[\texorpdfstring{\pa{ਛੋਟਾ} (chhoṭā)}{ਛੋਟਾ (chhoṭā)}]{\pa{ਛੋਟਾ} (chhoṭā) — small}

\label{lesson:PA-C79-chhota}



\begin{quote}
Before the new one: say the Punjabi for a snake, then the Punjabi for big.
\end{quote}

\subsection*{You'll want to know: \pa{ਛੋਟਾ}}
\textbf{\pa{ਛੋਟਾ}} — \emph{chhoṭā} — \textquotedblleft{}small\textquotedblright{}.

\begin{cousinweb}
\textbf{\pa{ਛੋਟਾ}} (\emph{chhoṭā}) is a tadbhava, a word inherited and reshaped by sound change. It comes from \emph{*choṭṭa}, \textquotedblleft{}small\textquotedblright{}, a form CDIAL reconstructs from its descendants (CDIAL 5071).
\end{cousinweb}

\begin{grammarlens}[title={What you've built}]
One more word for this chapter.
\end{grammarlens}

\subsection*{Your turn}
\begin{itemize}
\raggedright
  \item \emph{Say it:} \emph{chhoṭā}
  \item \emph{Say it:} \emph{chhoṭā}, once more
  \item \emph{Say it:} say \emph{chhoṭā}
  \item \emph{Recall:} say the Punjabi for a snake, then the Punjabi for big, then say \emph{chhoṭā} again
\end{itemize}

\begin{tcolorbox}[breakable,colback=teal!4,colframe=teal!35!black,title={Before you move on}]
What does \pa{ਛੋਟਾ} mean? (\textquotedblleft{}Small\textquotedblright{}.) Say it once more.
\end{tcolorbox}
\section[\texorpdfstring{\pa{ਨਵਾਂ} (navāṁ)}{ਨਵਾਂ (navāṁ)}]{\pa{ਨਵਾਂ} (navāṁ) — new}

\label{lesson:PA-C79-navam}



\begin{quote}
Before the new one: say the Punjabi for big, then the Punjabi for small.
\end{quote}

\subsection*{You'll want to know: \pa{ਨਵਾਂ}}
\textbf{\pa{ਨਵਾਂ}} — \emph{navāṁ} — \textquotedblleft{}new\textquotedblright{}.

\begin{cousinweb}
\textbf{\pa{ਨਵਾਂ}} (\emph{navāṁ}) is a tadbhava, a word inherited from Sanskrit and reshaped by sound change. It comes from \emph{nava}, \textquotedblleft{}new, young\textquotedblright{} (CDIAL 6983).
\end{cousinweb}

\begin{grammarlens}[title={What you've built}]
One more word for this chapter.
\end{grammarlens}

\subsection*{Your turn}
\begin{itemize}
\raggedright
  \item \emph{Say it:} \emph{navāṁ}
  \item \emph{Say it:} \emph{navāṁ}, once more
  \item \emph{Say it:} say \emph{navāṁ}
  \item \emph{Recall:} say the Punjabi for big, then the Punjabi for small, then say \emph{navāṁ} again
\end{itemize}

\begin{tcolorbox}[breakable,colback=teal!4,colframe=teal!35!black,title={Before you move on}]
What does \pa{ਨਵਾਂ} mean? (\textquotedblleft{}New\textquotedblright{}.) Say it once more.
\end{tcolorbox}
\section[\texorpdfstring{\pa{ਪੁਰਾਣਾ} (purāṇā)}{ਪੁਰਾਣਾ (purāṇā)}]{\pa{ਪੁਰਾਣਾ} (purāṇā) — old}

\label{lesson:PA-C79-purana}



\begin{quote}
Before the new one: say the Punjabi for small, then the Punjabi for new.
\end{quote}

\subsection*{You'll want to know: \pa{ਪੁਰਾਣਾ}}
\textbf{\pa{ਪੁਰਾਣਾ}} — \emph{purāṇā} — \textquotedblleft{}old\textquotedblright{}.

\begin{cousinweb}
\textbf{\pa{ਪੁਰਾਣਾ}} (\emph{purāṇā}) is inherited from Sanskrit \emph{purāṇa}, \textquotedblleft{}ancient\textquotedblright{}. CDIAL lists the matching Lahnda \emph{purāṇā} under it (CDIAL 8283).
\end{cousinweb}

\begin{grammarlens}[title={What you've built}]
One more word for this chapter.
\end{grammarlens}

\subsection*{Your turn}
\begin{itemize}
\raggedright
  \item \emph{Say it:} \emph{purāṇā}
  \item \emph{Say it:} \emph{purāṇā}, once more
  \item \emph{Say it:} say \emph{purāṇā}
  \item \emph{Recall:} say the Punjabi for small, then the Punjabi for new, then say \emph{purāṇā} again
\end{itemize}

\begin{tcolorbox}[breakable,colback=teal!4,colframe=teal!35!black,title={Before you move on}]
What does \pa{ਪੁਰਾਣਾ} mean? (\textquotedblleft{}Old\textquotedblright{}.) Say it once more.
\end{tcolorbox}
